% paper_full.tex — 全章节生产论文骨架（P8–P11，production 层）
% 定位：production 层的完整章节骨架（examples.demo-model 的对照约束：
%       本文件经受控执行与 Registry 登记后才可承载生产证据）。
% 纪律：
%   - 每个科学数字必须经 \MMResult{...}/\MMClaim{...}/\MMGiven{...} 等
%     绑定宏引用 Registry，禁止手抄数值；
%   - 每张图来自登记的生产 figure（\includegraphics 引用登记的产物文件）；
%   - 章节为结构示例：按题目增删小问，不做与题目无关的填充。
% 宏定义与 production/main.tex 保持一致（render-bindings 依此解析）。
\documentclass[UTF8,a4paper,12pt]{ctexart}
\usepackage[margin=2.5cm]{geometry}
\usepackage{amsmath,amssymb,booktabs,graphicx,hyperref}
\usepackage{listings}
% 占位渲染用 detokenize：占位 ID（如 __RESQ1__）含下划线，裸排印会触发
% 数学模式错误；detokenize 后在文本模式安全显示。
\newcommand{\MMFLOWUnresolved}[2]{\textbf{[MMFLOW-UNRESOLVED \detokenize{#1}: \detokenize{#2}]}}
\newcommand{\MMResult}[1]{\MMFLOWUnresolved{RESULT}{#1}}
\newcommand{\MMResultWithUnit}[1]{\MMFLOWUnresolved{RESULT-WITH-UNIT}{#1}}
\newcommand{\MMClaim}[2]{\MMFLOWUnresolved{CLAIM}{#1}: #2}
\newcommand{\MMGiven}[2]{\MMFLOWUnresolved{GIVEN}{#1}: #2}
\newcommand{\MMCitedValue}[2]{\MMFLOWUnresolved{CITED-VALUE}{#1}: #2}
\newcommand{\MMFormulaConstant}[2]{\MMFLOWUnresolved{FORMULA-CONSTANT}{#1}: #2}
\newcommand{\MMRuleValue}[2]{\MMFLOWUnresolved{RULE-VALUE}{#1}: #2}
\newcommand{\MMRequiredText}[1]{\textbf{[MMFLOW-UNRESOLVED TEXT: #1]}}
\lstset{basicstyle=\ttfamily\footnotesize,breaklines=true,
  keywordstyle=\color{blue!60!black},commentstyle=\color{green!45!black},
  showstringspaces=false,frame=single,numbers=left,numberstyle=\tiny}

\title{\MMRequiredText{论文标题（方法主线 + 问题对象，不写奖项暗示）}}
\author{}
\date{}
\begin{document}
\maketitle

\begin{abstract}
% 摘要六要素：背景动机 / 方法路线 / 关键结果（绑定数值）/ 验证证据 /
% 实际含义 / 限制。每段回答一个判断；证据未支持的内容减少主张而非补写。
\MMRequiredText{背景与方法总述}
\MMRequiredText{问题一方法与关键结果（\MMResult{...} 绑定）}
\MMRequiredText{问题二至问题N逐问段落}
\MMRequiredText{验证与灵敏度概括}
\MMRequiredText{实际含义与限制}
\label{abstractend}
\end{abstract}
% 摘要篇幅用 templates/publication/check_abstract_fill.py 复核（编译后运行）

\section{问题重述}
\subsection{问题背景}
\MMRequiredText{来自题面与一手规则的事实背景；外部信息用 \MMCitedValue{...} 绑定}
\subsection{问题提出}
\MMRequiredText{逐问重述：Q1/Q2/…/QN 各自要解决什么、产出什么形式}

\section{问题分析}
\MMRequiredText{逐问分析：输入—机制—输出链条、难点、与后续小问的依赖关系}

\section{模型假设}
\begin{enumerate}
  \item \MMRequiredText{假设1：内容 + \MMGiven{...} 依据 + 合理性论证}
  \item \MMRequiredText{假设2：…}
\end{enumerate}

\section{符号说明}
\begin{table}[htbp]
\centering
\begin{tabular}{cll}
\toprule
符号 & 含义 & 单位/来源 \\
\midrule
$x_i$ & \MMRequiredText{含义} & \MMRequiredText{单位；题面给定时 \MMGiven{...}} \\
\bottomrule
\end{tabular}
\end{table}

\section{问题一的建模与求解}
\subsection{建模推导}
\MMRequiredText{从假设到模型的完整推导；常数来源用 \MMFormulaConstant{...}；
规则性数字用 \MMRuleValue{...}}
\subsection{算法设计}
\MMRequiredText{算法步骤、复杂度与停止条件；伪代码可引用登记的 production 代码}
\subsection{结果与分析}
\begin{figure}[htbp]
\centering
% 文件必须是已登记 production figure 的产物；登记信息见 figure sidecar
\includegraphics[width=0.8\textwidth]{figures/model_result.png}
\caption{\MMRequiredText{图注四要素：问题、观察、单位与不确定性、解读}}
\end{figure}
\MMRequiredText{结果解读三层：数值事实（\MMResult{...}）→ 数学结构 → 决策含义}

\section{问题二至问题N的建模与求解}
% 每问重复"建模—算法—验证—结果"结构；跨问接口（共享变量/传递结果）显式
% 声明并与问题契约的依赖 DAG 一致
\MMRequiredText{逐问内容}

\section{模型验证与灵敏度分析}
\subsection{多方法交叉验证}
\MMRequiredText{题型适配验证结论（引用 validation_adapter_report 检查项）}
\subsection{误差分析}
\MMRequiredText{数值误差、离散化/收敛证据（可引用 convergence 扫描报告）}
\subsection{灵敏度分析}
\MMRequiredText{Tornado/Sobol/蒙特卡洛要点：最敏感因素、结论不变域}
\subsection{稳健性检验}
\MMRequiredText{退化场景、边界测试与最坏情况结论（\MMResult{...} 绑定）}

\section{模型评价与推广}
\subsection{优点}\MMRequiredText{与证据绑定的优点，不写空泛形容词}
\subsection{不足与改进}\MMRequiredText{诚实陈述 2--3 个具体局限与量化影响}
\subsection{推广}\MMRequiredText{适用边界内的推广路径；越界部分标为假设}

\section{结论}
\MMRequiredText{逐问最终答案汇总；只列证据充分的实质贡献}

\begin{thebibliography}{9}
\bibitem{ref1} \MMRequiredText{文献条目（GB/T 7714 或赛事指定格式）；
登记为 citation 后以 \MMCitedValue{...} 引用数值}
\end{thebibliography}

\appendix
\section{支撑材料清单}
\MMRequiredText{代码、配置、数据说明、复现说明与对应 artifact 说明}
\section{核心代码}
% 附录代码须与登记源文件逐字一致；用 templates/quality/
% appendix_code_closure_check.py 做闭包核验
\begin{lstlisting}[language=Python]
# 附录代码示例位置：替换为实际登记的生产代码文件内容
\end{lstlisting}
\end{document}
